\documentclass[11pt]{article}
\usepackage[margin=1in]{geometry}
\usepackage{xcolor}
\usepackage{listings}
\usepackage{booktabs}

\lstset{
  basicstyle=\ttfamily\small,
  keywordstyle=\bfseries\color{blue!60!black},
  commentstyle=\itshape\color{gray!70!black},
  stringstyle=\color{green!35!black},
  numbers=left,
  numberstyle=\tiny,
  stepnumber=1,
  numbersep=6pt,
  frame=single,
  breaklines=true,
  columns=fullflexible
}

\title{Source Listing Stress Case}
\author{Test Corpus}
\date{}

\begin{document}
\maketitle

\begin{abstract}
This fixture tests monospaced inline text and a displayed source listing using
the \texttt{listings} package. The example resembles a reproducibility appendix
in a computational social science paper.
\end{abstract}

\section{Pipeline Description}
\label{sec:pipeline}

The analysis pipeline begins with the function \texttt{clean\_panel()} and ends
with a short diagnostic table. Inline code should remain monospaced when the
document is converted, including identifiers such as \texttt{site\_id} and
\texttt{event\_time}.

\begin{lstlisting}[language=Python,caption={Minimal synthetic cleaning step},label={lst:cleaning}]
def clean_panel(records):
    """Return observations with complete identifiers."""
    cleaned = []
    for row in records:
        if row["site_id"] and row["event_time"] is not None:
            row["post"] = int(row["event_time"] >= 0)
            cleaned.append(row)
    return cleaned
\end{lstlisting}

Listing~\ref{lst:cleaning} includes indentation, comments, strings, line
numbers, and a caption. These elements commonly reveal whether a converter treats
source listings as formatted code blocks or as ordinary paragraphs.

\section{Diagnostic Output}
\label{sec:diagnostic}

\begin{table}[htbp]
\centering
\caption{Synthetic code-path diagnostics}
\label{tab:diagnostics}
\begin{tabular}{lcc}
\toprule
Check & Rows affected & Status \\
\midrule
Identifier present & 240 & Pass \\
Event time present & 233 & Pass \\
Post indicator created & 233 & Pass \\
\bottomrule
\end{tabular}
\end{table}

The table confirms that the listing is embedded in an otherwise conventional
academic document with cross-references and prose.

\end{document}
